\section{USAMO 2001}
        \begin{enumerate}
        	\item (\textit{USAMO 2001}) Each of eight boxes contains six balls. Each ball has been colored with one of $n$ colors, such that no two balls in the same box are the same color, and no two colors occur together in more than one box. Determine, with justification, the smallest integer $n$ for which this is possible.


 



\textbf{Solution 1}

We claim that $n=23$ is the minimum. Consider the following construction (replacing colors with numbers) which fulfills this:
 \[\left[ \begin{array}{cccccccc} 1 & 1 & 1 & 2 & 3 & 4 & 5 & 6 \\ 2 & 7 & 12 & 7 & 8 & 9 & 10 & 11 \\ 3 & 8 & 13 & 12 & 13 & 14 & 15 & 16 \\ 4 & 9 & 14 & 17 & 17 & 17 & 18 & 19 \\ 5 & 10 & 15 & 18 & 20 & 22 & 20 & 21 \\ 6 & 11 & 16 & 19 & 21 & 23 & 22 & 23 \end{array} \right]\]Suppose a configuration exists with $n \le 22$. 

 Suppose a ball appears $5$ or more times. Then the remaining balls of the $5$ boxes must be distinct, so that there are at least $n \ge 5 \cdot 5 + 1 = 26$ balls, contradiction. If a ball appears $4$ or more times, the remaining balls of the $4$ boxes must be distinct, leading to $5 \cdot 4 + 1 = 21$ balls. The fifth box can contain at most four balls from the previous boxes, and then the remaining two balls must be distinct, leading to $n \ge 2 + 21 = 23$, contradiction.

 However, by the \href{https://artofproblemsolving.com/wiki/index.php?title=Pigeonhole\_Principle}{Pigeonhole Principle}, at least one ball must appear $3$ times. \href{https://artofproblemsolving.com/wiki/index.php?title=Without\_loss\_of\_generality}{Without loss of generality} suppose that $1$ appears three times, and let the boxes that contain these have balls with colors $\{1,2,3,4,5,6\},\{1,7,8,9,10,11\},\{1,12,13,14,15,16\}$. Each of the remaining five boxes can have at most $3$ balls from each of these boxes. Thus, each of the remaining five boxes must have $3$ additional balls $> 16$. Thus, it is necessary that we use $\le 22 - 16 = 6$ balls to fill a $3 \times 5$ grid by the same rules. 

 Again, no balls may appear $\ge 4$ times, but by Pigeonhole, one ball must appear $3$ times. \href{https://artofproblemsolving.com/wiki/index.php?title=Without\_loss\_of\_generality}{Without loss of generality}, let this ball have color $17$; then the three boxes containing $17$ must have at least $2 \cdot 3 + 1 = 7$ balls, contradiction.

 Therefore, $n = 23$ is the minimum.

 



\textbf{Solution 2}

Similar to the above solution, no color can appear 4 times or more.

 We can use PIE. Let $S_1,S_2,\cdots ,S_8$ be the colors in each of the $8$ boxes. Then $n=|\cup_{i=1}^8S_i|$. By PIE we know that
 \[|\cup_{i=1}^8S_i|=\sum_{i=1}^8 |S_i|-\sum_{1\leq i<j\leq 8}|S_i\cap S_j|+\sum_{1\leq i<j<k\leq 8}|S_i\cap S_j\cap S_k|-\sum_{1\leq i<j<k<l\leq 8}|S_i\cap S_j\cap S_k\cap S_l|\cdots.\]Note however that no color can appear 4 times or more, so all the items after $\sum_{1\leq i<j<k\leq 8}|S_i\cap S_j\cap S_k|$ must be equal to 0, so
 \[|\cup_{i=1}^8S_i|=\sum_{i=1}^8 |S_i|-\sum_{1\leq i<j\leq 8}|S_i\cap S_j|+\sum_{1\leq i<j<k\leq 8}|S_i\cap S_j\cap S_k|.\]Now note that for each $i$, $|S_i|=6,$ so $\sum_{i=1}^8 |S_i|=48.$ Also no pair of colors appears twice, so for any $i,j, |S_i\cap S_j|\leq 1,$ and so $\sum_{1\leq i<j\leq 8}|S_i\cap S_j|\leq \binom82=28.$ Thus, $n\geq 48-28+\sum_{1\leq i<j<k\leq 8}|S_i\cap S_j\cap S_k|.$ Given that there are $n$ colors, to minimize the number of colors in 3 boxes means each color is in either 2 or 3 boxes, and since there are 48 total spots, this means there are $48-2n$ colors in 3 boxes and $3n-48$ colors in 2. Thus $n\geq 20+48-2n,$ so $3n\geq 68$, and $n\geq23$ and we are done.

 



\textbf{Solution 3}

Let $m_{i,j}$ be the number of balls which are the same color as the $j^{\text{th}}$ ball in box $i$ (including that ball). For a fixed box $i$, $1\leq i\leq 8$, consider the sums
 \[S_i = \sum_{j=1}^6 m_{i,j}\quad\text{and}\quad s_i = \sum_{j=1}^6 \frac{1}{m_{i,j}}.\]For each fixed $i$, since no pair of colors is repeated, each of the remaining seven boxes can contribute at most one ball to $S_i$. Thus, $S_i\leq 13$. It follows by the convexity of $f(x) = 1/x$ that $s_i$ is minimized when one of the $m_{i,j}$ is equal to 3 and the other five equal to 2. Hence $s_i\geq 17/6$. Note that
 \[n = \sum_{i=1}^8 \sum_{j=1}^6 \frac{1}{m_{i,j}}\geq 8\cdot\frac{17}{6} = \frac{68}{3}.\]Hence there must be 23 colors. The construction for $n = 23$ is above.

 


	\item (\textit{USAMO 2001}) Let $ABC$ be a \href{https://artofproblemsolving.com/wiki/index.php?title=Triangle}{triangle} and let $\omega$ be its \href{https://artofproblemsolving.com/wiki/index.php?title=Incircle}{incircle}. Denote by $D_1$ and $E_1$ the points where $\omega$ is tangent to sides $BC$ and $AC$, respectively. Denote by $D_2$ and $E_2$ the points on sides $BC$ and $AC$, respectively, such that $CD_2 = BD_1$ and $CE_2 = AE_1$, and denote by $P$ the point of intersection of segments $AD_2$ and $BE_2$. Circle $\omega$ intersects segment $AD_2$ at two points, the closer of which to the vertex $A$ is denoted by $Q$. Prove that $AQ = D_2P$.


 



\textbf{Solution 1}

It is well known that the excircle opposite $A$ is tangent to $\overline{BC}$ at the point $D_2$. (Proof: let the points of tangency of the excircle with the lines $BC, AB, AC$ be $D_3, F,G$ respectively. Then $AB+BD_3=AB + BF=AF = AG = AC + CG=AC + CD_3$. It follows that $2CD_3 = AB + BC - AC$, and $CD_3 = s-b = BD_1 = CD_2$, so $D_3 \equiv D_2$.)

 Now consider the homothety that carries the incircle of $\triangle ABC$ to its excircle. The homothety also carries $Q$ to $D_2$ (since $A,Q,D_2$ are collinear), and carries the tangency points $E_1$ to $G$. It follows that $\frac{AQ}{QD_2} = \frac{AE_1}{E_1G} = \frac{s-a}{E_1C + CD_2} = \frac{s-a}{CD_1 + BD_1} = \frac{s-a}{a}$.

 
                    \begin{figure}[H]
                        \centering
                        \includegraphics[width=0.5\textwidth]{Solution/USAMO_I_2001_2_Solution_Figure_1.png}
                    \end{figure}
                    

 By Menelaus' Theorem on $\triangle ACD_2$ with segment $\overline{BE_2}$, it follows that $\frac{CE_2}{E_2A} \cdot \frac{AP}{PD_2} \cdot \frac{BD_2}{BC} = 1 \Longrightarrow \frac{AP}{PD_2} = \frac{(b - (s-a)) \cdot a}{(a-(s-b)) \cdot AE_1} = \frac{a}{s-a}$. It easily follows that $AQ = D_2P$. $\blacksquare$

 



\textbf{Solution 2}

The key observation is the following lemma.

 Lemma: Segment $D_1Q$ is a diameter of circle $\omega$.

 Proof: Let $I$ be the center of circle $\omega$, i.e., $I$ is the incenter of triangle $ABC$. Extend segment $D_1I$ through $I$ to intersect circle $\omega$ again at $Q'$, and extend segment $AQ'$ through $Q'$ to intersect segment $BC$ at $D'$. We show that $D_2 = D'$, which in turn implies that $Q = Q'$, that is, $D_1Q$ is a diameter of $\omega$.

 Let $l$ be the line tangent to circle $\omega$ at $Q'$, and let $l$ intersect the segments $AB$ and $AC$ at $B_1$ and $C_1$, respectively. Then $\omega$ is an excircle of triangle $AB_1C_1$. Let $\mathbf{H}_1$ denote the dilation with its center at $A$ and ratio $AD'/AQ'$. Since $l\perp D_1Q'$ and $BC\perp D_1Q'$, $l\parallel BC$. Hence $AB/AB_1 = AC/AC_1 = AD'/AQ'$. Thus $\mathbf{H}_1(Q') = D'$, $\mathbf{H}_1(B_1) = B$, and $\mathbf{H}_1(C_1) = C$. It also follows that an excircle $\Omega$ of triangle $ABC$ is tangent to the side $BC$ at $D'$.

 It is well known that 
 \[CD_1 = \frac{1}{2}(BC + CA - AB).\] We compute $BD'$. Let $X$ and $Y$ denote the points of tangency of circle $\Omega$ with rays $AB$ and $AC$, respectively. Then by equal tangents, $AX = AY$, $BD' = BX$, and $D'C = YC$. Hence 
 \[AX = AY = \frac{1}{2}(AX + AY) = \frac{1}{2}(AB + BX + YC + CA) = \frac{1}{2}(AB + BC + CA).\] It follows that 
 \[BD' = BX = AX - AB = \frac{1}{2}(BC + CA - AB).\] Combining these two equations yields $BD' = CD_1$. Thus 
 \[BD_2 = BD_1 - D_2D_1 = D_2C - D_2D_1 = CD_1 = BD',\] that is, $D' = D_2$, as desired. $\blacksquare$

 Now we prove our main result. Let $M_1$ and $M_2$ be the respective midpoints of segments $BC$ and $CA$. Then $M_1$ is also the midpoint of segment $D_1D_2$, from which it follows that $IM_1$ is the midline of triangle $D_1QD_2$. Hence 
 \[QD_2 = 2IM_1\] and $AD_2\parallel M_1I$. Similarly, we can prove that $BE_2\parallel M_2I$.

 2001usamo2-2.pngLet $G$ be the centroid of triangle $ABC$. Thus segments $AM_1$ and $BM_2$ intersect at $G$. Define transformation $\mathbf{H}_2$ as the dilation with its center at $G$ and ratio $-1/2$. Then $\mathbf{H}_2(A) = M_1$ and $\mathbf{H}_2(B) = M_2$. Under the dilation, parallel lines go to parallel lines and the intersection of two lines goes to the intersection of their images. Since $AD_2\parallel M_1I$ and $BE_2\parallel M_2I$, $\mathbf{H}_2$ maps lines $AD_2$ and $BE_2$ to lines $M_1I$ and $M_2I$, respectively. It also follows that $\mathbf{H}_2(I) = P$ and 
 \[\frac{IM_1}{AP} = \frac{GM_1}{AG} = \frac{1}{2}\] or 
 \[AP = 2IM_1.\] This yields 
 \[AQ = AP - QP = 2IM_1 - QP = QD_2 - QP = PD_2,\] as desired.

 Note: We used directed lengths in our calculations to avoid possible complications caused by the different shapes of triangle $ABC$.

 



\textbf{Solution 3}

Here is a rather nice solution using barycentric coordinates:

 Let $A$ be $(1,0,0)$, $B$ be $(0,1,0)$, and $C$ be $(0,0,1)$. Let the side lengths of the triangle be $a,b,c$ and the semi-perimeter $s$.

 Now, 
 \[CD_1=s-c, BD_1=s-b, AE_1=s-a, CE_1=s-c.\] Thus, 
 \[CD_2=s-b, BD_2=s-c, AE_2=s-c, CE_2=s-a.\]

 Therefore, $D_2=(0:s-b:s-c)$ and $E_2=(s-a:0:s-c).$ Clearly then, the non-normalized coordinates of $P=(s-a:s-b:s-c).$

 Normalizing, we have that 
 \[D_2=\left(0,\frac{s-b}{a},\frac{s-c}{a}\right), E_2=\left(\frac{s-a}{b},0,\frac{s-c}{b}\right), P=\left(\frac{s-a}{s},\frac{s-b}{s},\frac{s-c}{s}\right).\]

 Now, we find the point $Q'$ inside the triangle on the line $AD_2$ such that $AQ'=D_2P$. It is then sufficient to show that this point lies on the incircle.

 $P$ is the fraction $\frac{s-a}{s}$ of the way "up" the line segment from $D_2$ to $A$. Thus, we are looking for the point that is $\frac{s-a}{s}$ of the way "down" the line segment from $A$ to $D_2$, or, the fraction $1-\frac{s-a}{s}$ of the way "up".

 Thus, $Q'$ has normalized $x$-coordinate $1-\frac{s-a}{s}=\frac{a}{s}$.

 As the line $AD_2$ has equation $(s-c)y=(s-b)z$, it can easily be found that 
 \[Q'=\left(\frac{a^2}{as}, \frac{(s-a)(s-b)}{as}, \frac{(s-a)(s-c)}{as}\right)=(a^2:(s-a)(s-b):(s-a)(s-c)).\]

 Recalling that the equation of the incircle is 
 \[a^2yz+b^2xz+c^2xy+(x+y+z)[(s-a)^2x+(s-b)^2y+(s-c)^2z]=0.\] We must show that this equation is true for $Q'$'s values of $x,y,z$.

 Plugging in our values, this means showing that
 \[a^2(s-a)^2(s-b)(s-c)+a^2(s-a)[b^2(s-c)+c^2(s-b)]+as[a^2(s-a)^2+(s-a)(s-b)^3+(s-a)(s-c)^3]=0.\]Dividing by $a(s-a)$, this is just
 \[a(s-a)(s-b)(s-c)+a[b^2(s-c)+c^2(s-b)]+a^2s(s-a)+s(s-b)^3+s(s-c)^3=0.\]

 Plugging in the value of $s:$
 \[\frac{a(-a+b+c)(a-b+c)(a+b-c)}{8}+\frac{ab^2(a+b-c)}{2}+\frac{ac^2(a-b+c)}{2}+\frac{a^2(a+b+c)(-a+b+c)}{4}+\frac{(a+b+c)(a-b+c)^3}{16}+\frac{(a+b+c)(a+b-c)^3}{16}=0.\]
 \[2a(-a+b+c)(a-b+c)(a+b-c)+8ab^2(a+b-c)+8ac^2(a-b+c)+4a^2(a+b+c)(-a+b+c)+(a+b+c)(a-b+c)^3+(a+b+c)(a+b-c)^3=0\]
 \[2a[(-a+b+c)(a-b+c)(a+b-c)+4b^2(a+b-c)+4c^2(a-b+c)]+(a+b+c)[4a^2(-a+b+c)+(a-b+c)^3+(a+b-c)^3]=0\]The first bracket is just
 \[-a^3-b^3-c^3+a^2b+a^2c+b^2c+b^2a+c^2a+c^2b-2abc+4ab^2+4b^3-4b^2c+4ac^2-4bc^2+4c^3\]
 \[=-a^3+3b^3+3c^3+a^2b+a^2c-3b^2c+5b^2a+5c^2a-3c^2b-2abc\]and the second bracket is 
 \[-2a^3+4a^2b+4a^2c+6ab^2+6ac^2-12abc.\]Dividing everything by $2a$ gives
 \[-a^3+3b^3+3c^3+a^2b+a^2c-3b^2c+5b^2a+5c^2a-3c^2b-2abc+(a+b+c)(-a^2+2ab+2ac+3b^2+3c^2-6bc),\] which is $0$, as desired.

 As $Q'$ lies on the incircle and $AD_2$, $Q'=Q$, and our proof is complete.

 


 



\textbf{Solution 4}

We again use Barycentric coordinates. As before, let $A$ be $(1,0,0)$, $B$ be $(0,1,0)$, and $C$ be $(0,0,1)$. Also

 
 \[D_{1}=\left( 0,\frac{s-c}{a},\frac{s-b}{a} \right), D_2=\left(0,\frac{s-b}{a},\frac{s-c}{a}\right), E_2=\left(\frac{s-a}{b},0,\frac{s-c}{b}\right), P=\left(\frac{s-a}{s},\frac{s-b}{s},\frac{s-c}{s}\right).\]

 Now, consider a point $Q'$ for which $AQ'=PD_{2}$. Then working component-vise, we get $Q'+P=A+D_{2}$ from which we can easily get the coordinates of $Q'$ as;

 
 \[Q'=\left(\frac{a}{s},\frac{(s-a)(s-b)}{sa},\frac{(s-a)(s-c)}{sa}\right)\]

 It suffices to show that $Q'=Q$.

 Let $I=\left(\frac{a}{2s},\frac{b}{2s},\frac{c}{2s}\right)$  be the incenter of triangle $ABC$. We claim that $I$ is the midpoint of $Q'D_{1}$. Indeed,

 
 \[\frac{1}{2}\left(0+\frac{a}{s}\right)=\frac{a}{2s}\]

 
 \[\frac{1}{2}\left(\frac{(s-a)(s-b)}{sa}+\frac{s-c}{a}\right)=\frac{(s-a)(s-b)+s(s-c)}{2sa}=\frac{ab}{2sa}=\frac{b}{2s}\]

 
 \[\frac{1}{2}\left(\frac{(s-a)(s-c)}{sa}+\frac{s-b}{a}\right)=\frac{c}{2s}\]

 Hence the claim has been proved.

 Since $I$ is the center of $\omega$ and the midpoint of  $Q'D_{1}$, thus $Q'$ is the point diametrically opposite to $D_{1}$, and hence it lies on $\omega$ and closer to $A$ off the 2 points. Thus $Q=Q'$ as desired.

 


	\item (\textit{USAMO 2001}) Let $a, b, c \geq 0$ and satisfy


 $a^2 + b^2 + c^2 + abc = 4.$

 Show that


 $0\leq ab + bc + ca - abc \leq 2.$

 



\textbf{Solution 1}

First we prove the lower bound.

 Note that we cannot have $a, b, c$ all greater than 1.Therefore, suppose $a \le 1$.Then
 \[ab + bc + ca - abc = a(b + c) + bc(1-a) \ge 0.\]Note that, by the \href{https://artofproblemsolving.com/wiki/index.php?title=Pigeonhole\_Principle}{Pigeonhole Principle}, at least two of $a,b,c$ are either both greater than or less than $1$. \href{https://artofproblemsolving.com/wiki/index.php?title=Without\_loss\_of\_generality}{Without loss of generality}, let them be $b$ and $c$. Therefore, $(b-1)(c-1)\ge 0$. From the given equation, we can express $a$ in terms of $b$ and $c$ as
 \[a=\frac{\sqrt{(4-b^2)(4-c^2)}-bc}{2}\]Thus,
 \[ab + bc + ca - abc = -a (b-1)(c-1)+a+bc \le a+bc = \frac{\sqrt{(4-b^2)(4-c^2)} + bc}{2}\]

 From the \href{https://artofproblemsolving.com/wiki/index.php?title=Cauchy-Schwarz\_Inequality}{Cauchy-Schwarz Inequality},
 \[\frac{\sqrt{(4-b^2)(4-c^2)} + bc}{2} \le \frac{\sqrt{(4-b^2+b^2)(4-c^2+c^2)} }{2} = 2.\]

 This completes the proof.

 



\textbf{Solution 2}

The proof for the lower bound is the same as in the first solution.

 Now we prove the upper bound. Let us note that at least two of the three numbers $a$, $b$, and $c$ are both greater than or equal to 1 or less than or equal to 1. Without loss of generality, we assume that the numbers with this property are $b$ and $c$. Then we have
 \[(1 - b)(1 - c)\geq 0.\]The given equality $a^2 + b^2 + c^2 + abc = 4$ and the inequality $b^2 + c^2\geq 2bc$ imply that
 \[a^2 + 2bc + abc\leq 4,\]or
 \[bc(2 + a)\leq 4 - a^2.\]Dividing both sides of the last inequality by $2 + a$ yields
 \[bc\leq 2 - a.\]Thus,
 \[ab + bc + ca - abc\leq ab + 2 - a + ac(1 - b) = 2 - a(1 + bc - b - c) = 2 - a(1 - b)(1 - c)\leq 2,\]as desired.

 The last equality holds if and only if $b = c$ and $a(1 - b)(1 - c) = 0$. Hence equality for the upper bound holds if and only if $(a,b,c)$ is one of the triples $(1,1,1)$, $(0,\sqrt{2},\sqrt{2})$, $(\sqrt{2},0,\sqrt{2})$, and $(\sqrt{2},\sqrt{2},0)$. Equality for the lower bound holds if and only if $(a,b,c)$ is one of the triples $(2,0,0)$, $(0,2,0)$ and $(0,0,2)$.

 



\textbf{Solution 3}

The proof for the lower bound is the same as in the first solution.

 Now we prove the upper bound. It is clear that $a,b,c\leq 2$. If $abc = 0$, then the result is trivial. Suppose that $a,b,c > 0$. Solving for $a$ yields
 \[a = \frac{-bc + \sqrt{b^2c^2 - 4(b^2 + c^2 - 4)}}{2} = \frac{-bc + \sqrt{(4 - b^2)(4 - c^2)}}{2}.\]This asks for the trigonometric substitution $b = 2\sin u$ and $c = 2\sin v$, where $0^\circ < u,v < 90^\circ$. Then
 \[a = 2(-\sin u\sin v + \cos u\cos v) = 2\cos (u + v),\]and if we set $u = B/2$ and $v = C/2$, then $a = 2\sin (A/2)$, $b = 2\sin (B/2)$, and $c = \sin (C/2)$, where $A$, $B$, and $C$ are the angles of a triangle. We have
 \begin{align*} ab &= 4\sin\frac{A}{2}\sin\frac{B}{2} \\ &= 2\sqrt{\sin A\tan\frac{A}{2}\sin B\tan\frac{B}{2}} = 2\sqrt{\sin A\tan\frac{B}{2}\sin B\tan\frac{A}{2}} \\ &\leq \sin A\tan\frac{B}{2} + \sin B\tan\frac{A}{2} \\ &= \sin A\cot\frac{A + C}{2} + \sin B\cot\frac{B + C}{2}, \end{align*}
where the inequality step follows from AM-GM. Likewise,
 \begin{align*} bc &\leq \sin B\cot\frac{B + A}{2} + \sin C\cot\frac{C + A}{2}, \\ ca &\leq \sin A\cot\frac{A + B}{2} + \sin C\cot\frac{C + B}{2}. \end{align*}
Therefore
 \begin{align*} ab + bc + ca &\leq (\sin A + \sin B)\cot\frac{A + B}{2} + (\sin B + \sin C)\cot\frac{B + C}{2} + (\sin C + \sin A)\cot\frac{C + A}{2} \\ &= 2\left(\cos\frac{A - B}{2}\cos\frac{A + B}{2} + \cos\frac{B - C}{2}\cos\frac{B + C}{2} + \cos\frac{C - A}{2}\cos\frac{C + A}{2} \right)\\ &= 2(\cos A + \cos B + \cos C) \\ &= 6 - 4\left(\sin^2\frac{A}{2} + \sin^2\frac{B}{2} + \sin^2\frac{C}{2}\right) \\ &= 6 - (a^2 + b^2 + c^2) = 2 + abc, \end{align*}
as desired.

 


	\item (\textit{USAMO 2001}) Let $P$ be a point in the plane of triangle $ABC$ such that the segments $PA$, $PB$, and $PC$ are the sides of an obtuse triangle. Assume that in this triangle the obtuse angle opposes the side congruent to $PA$. Prove that $\angle BAC$ is acute.


 



\textbf{Solution 1}

We know that $PB^2+PC^2 < PA^2$ and we wish to prove that $AB^2 + AC^2 > BC^2$.It would be sufficient to prove that
 \[PB^2+PC^2+AB^2+AC^2 \geq PA^2 + BC^2.\]Set $A(0,0)$, $B(1,0)$, $C(x,y)$, $P(p,q)$.Then, we wish to show

 
 \[(p-1)^2 + q^2 + (p-x)^2 + (q-y)^2 + 1 + x^2 + y^2 \geq p^2 + q^2 + (x-1)^2 + y^2\]
 \[2p^2 + 2q^2 + 2x^2 + 2y^2 - 2p - 2px - 2qy + 2 \geq p^2 + q^2 + x^2 + y^2 - 2x + 1\]
 \[p^2 + q^2 + x^2 + y^2 + 2x - 2p - 2px - 2qy + 1 \geq 0\]
 \[(x-p)^2 + (q-y)^2 + 2(x-p) + 1 \geq 0\] 
 \[(x-p+1)^2 + (q-y)^2 \geq 0,\]

 which is true by the trivial inequality.

 



\textbf{Solution 2}

Let $A$ be the origin. For a point $Q$, denote by $q$ the vector $\overrightarrow{AQ}$, and denote by $|q|$ the length of $q$. The given conditions may be written as
 \[|p - b|^2 + |p - c|^2 < |p|^2,\]or
 \[p\cdot p + b\cdot b + c\cdot c - 2p\cdot b - 2p\cdot c < 0.\]Adding $2b\cdot c$ on both sides of the last inequality gives
 \[|p - b - c|^2 < 2b\cdot c.\]Since the left-hand side of the last inequality is nonnegative, the right-hand side is positive. Hence
 \[\cos\angle BAC = \frac{b\cdot c}{|b||c|} > 0,\]that is, $\angle BAC$ is acute.

 



\textbf{Solution 3}

For the sake of contradiction, let's assume to the contrary that $\angle BAC$. Let $AB = c$, $BC = a$, and $CA = b$. Then $a^2\geq b^2 + c^2$. We claim that the quadrilateral $ABPC$ is convex. Now applying the generalized Ptolemy's Theorem to the convex quadrilateral $ABPC$ yields
 \[a\cdot PA\leq b\cdot PB + c\cdot PC\leq\sqrt{b^2 + c^2}\sqrt{PB^2 + PC^2}\leq a\sqrt{PB^2 + PC^2},\]where the second inequality is by Cauchy-Schwarz. This implies $PA^2\leq PB^2 + PC^2$, in contradiction with the facts that $PA$, $PB$, and $PC$ are the sides of an obtuse triangle and $PA > \max\{PB, PC\}$.

 We present two arguments to prove our claim.

 \textit{First argument}: Without loss of generality, we may assume that $A$, $B$, and $C$ are in counterclockwise order. Let lines $l_1$ and $l_2$ be the perpendicular bisectors of segments $AB$ and $AC$, respectively. Then $l_1$ and $l_2$ meet at $O$, the circumcenter of triangle $ABC$. Lines $l_1$ and $l_2$ cut the plane into four regions and $A$ is in the interior of one of these regions. Since $PA > PB$ and $PA > PC$, $P$ must be in the interior of the region that opposes $A$. Since $\angle BAC$ is not acute, ray $AC$ does not meet $l_1$ and ray $AB$ does not meet $l_2$. Hence $B$ and $C$ must lie in the interiors of the regions adjacent to $A$. Let $\mathcal{R}_X$ denote the region containing $X$. Then $\mathcal{R}_A$, $\mathcal{R}_B$, $\mathcal{R}_P$, and $\mathcal{R}_C$ are the four regions in counterclockwise order. Since $\angle BAC\geq 90^\circ$, either $O$ is on side $BC$ or $O$ and $A$ are on opposite sides of line $BC$. In either case $P$ and $A$ are on opposite sides of line $BC$. Also, since ray $AB$ does not meet $l_2$ and ray $AC$ does not meet $l_1$, it follows that $\mathcal{R}_P$ is entirely in the interior of $\angle BAC$. Hence $B$ and $C$ are on opposite sides of $AP$. Therefore $ABPC$ is convex.

 
                    \begin{figure}[H]
                        \centering
                        \includegraphics[width=0.5\textwidth]{Solution/USAMO_I_2001_4_Solution_Figure_1.png}
                    \end{figure}
                    

 \textit{Second argument}: Since $PA > PB$ and $PA > PC$, $A$ cannot be inside or on the sides of triangle $PBC$. Since $PA > PB$, we have $\angle ABP > \angle BAP$ and hence $\angle BAC\geq 90^\circ > \angle BAP$. Hence $C$ cannot be inside or on the sides of triangle $BAP$. Symmetrically, $B$ cannot be inside or on the sides of triangle $CAP$. Finally, since $\angle ABP > \angle BAP$ and $\angle ACP > \angle CAP$, we have
 \[\angle ABP + \angle ACP > \angle BAC\geq 90^\circ\geq\angle ABC + \angle ACB.\]Therefore $P$ cannot be inside or on the sides of triangle $ABC$. Since this covers all four cases, $ABPC$ is convex.

 



\textbf{Solution 4}

Let $P$ be the origin in vector space, and let $a, b, c$ denote the position vectors of $A, B, C$ respectively. Then the obtuse triangle condition, $PA^2 > PB^2 + PC^2$, becomes $a^2 > b^2 + c^2$ using the fact that the square of a vector (the dot product of itself and itself) is the square of its magnitude. Now, notice that to prove $\angle{BAC}$ is acute, it suffices to show that $(a - b)(a - c) > 0$, or $a^2 - ab - ac + bc > 0$. But this follows from the observation that
 \[(-a + b + c)^2 \ge 0,\]which leads to
 \[2a^2 - 2ab - 2ac + 2bc > a^2 + b^2 + c^2 - 2ab - 2ac + 2bc \ge 0\]and therefore our desired conclusion.

 



\textbf{Solution 5}

Let $M, N$ be midpoints of $AP$ and $BC$, respectively. For the points $A, B, P, C$; let's apply Euler's quadrilateral formula,
 \[AB^2 + BP^2 + PC^2 + CA^2 = AP^2 + BC^2 + 4MN^2 \geq AP^2 + BC^2 .\]Given that $AP^2 > BP^2 + PC^2$. Thus,
 \[AB^2 + AC^2 > BC^2 .\] and we get $\angle BAC$ is acute.

 (Lokman GÖKÇE)

 



\textbf{Solution 6}

Without loss of generality, assume that in a Cartesian coordinate system, $A$ is at the point $(0, 0)$ and $C$ is at the point $(1,0)$. Let $B$ be at the point $(b_x,b_y)$ and $P$ be at the point $(p_x,p_y)$. Without loss of generality, also assume that $b_y>0$.

 Now, assume for contradiction that $\angle BAC$ is not acute. Since $PA$, $PB$, and $PC$ are the sides of an obtuse triangle, with $PA$ the longest side, it follows that $PA^2>PB^2+PC^2$, implying that $p_x^2+p_y^2>(p_x-b_x)^2+(p_y-b_y)^2+(p_x-1)^2+p_y^2$. This inequality simplifies to $b_x^2-2p_x b_x+b_y^2-2p_y b_y+p_x^2-2p_x+1+p_y^2<0$. Note that since $p_x^2-2p_x+1$ and $b_y^2-2p_y b_y+p_y^2$ are both perfect squares, all terms of this inequality except for $-2p_x b_x$ are already guaranteed to be nonnegative.

 If $p_x<0$, then $P$ would be closer to $A$ than to $C$, but since $PA^2=PB^2+PC^2$, this is not possible. Therefore, $p_x \geq 0$. Since $\angle BAC$ not being acute implies that $b_x \leq 0$, it follows that $-2p_x b_x \geq 0$. But now since all terms of $b_x^2-2p_x b_x+b_y^2-2p_y b_y+p_x^2-2p_x+1+p_y^2<0$ are guaranteed to be nonnegative, this entire expression cannot be negative, leading to a contradiction. Therefore, $\angle BAC$ is acute.

 



\textbf{Solution 7}

To maximize $\angle BAC$, we try to maximize $\angle PAB$ and $\angle PAC$, which occurs when $PB = PA \sin \angle PAB$ and $PC = PA \sin \angle PAC$. Let $\angle PAB = \theta_1$ and let $\angle PAC = \theta_2$. 

 Since $PC, PB$, and $PA$ form an obtuse triangle, $PC^2 + PB^2 < PA^2$. 

 We now have that $\sin^2(\theta_1) + \sin^2(\theta_2) = \frac{PB^2+PC^2}{PA^2} < 1$. The premise is now obvious; $\sin^2(\theta_2) < \cos^2(\theta_1)$, which means that $\theta_2 < 90 - \theta_1$, so we're done.

 


	\item (\textit{USAMO 2001}) Let $S$ be a set of integers (not necessarily positive) such that


 (a) there exist $a,b \in S$ with $\gcd(a,b) = \gcd(a - 2,b - 2) = 1$;


 (b) if $x$ and $y$ are elements of $S$ (possibly equal), then $x^2 - y$ also belongs to $S$. 


 Prove that $S$ is the set of all integers.


 



\textbf{Solution}

In the solution below we use the expression $S$ \textit{is stable under} $x\mapsto f(x)$ to mean that if $x$ belongs to $S$ then $f(x)$ also belongs to $S$. If $c,d\in S$, then by (B), $S$ is stable under $x\mapsto c^2 - x$ and $x\mapsto d^2 - x$, hence stable under $x\mapsto c^2 - (d^2 - x) = x + (c^2 - d^2)$. Similarly $S$ is stable under $x\mapsto x + (d^2 - c^2)$. Hence $S$ is stable under $x\mapsto x + n$ and $x\mapsto x - n$ whenever $n$ is an integer linear combination of numbers of the form $c^2 - d^2$ with $c,d\in S$. In particular, this holds for $n = m$, where $m = \gcd\{c^2 - d^2 : c,d\in S\}$.

 Since $S\neq \emptyset$ by (A), it suffices to prove that $m = 1$. For the sake of contradiction, assume that $m\neq 1$. Let $p$ be a prime dividing $m$. Then $c^2 - d^2\equiv 0\pmod{p}$ for all $c,d\in S$. In other words, for each $c,d\in S$, either $d\equiv c\pmod{p}$ or $d\equiv -c\pmod{p}$. Given $c\in S$, $c^2 - c\in S$ by (B), so $c^2 - c\equiv c\pmod{p}$ or $c^2 - c\equiv -c\pmod{p}$. Hence
 \[\text{For each }c\in S\text{, either }c\equiv 0\pmod{p}\text{ or }c\equiv 2\pmod{p}.\qquad\qquad (*)\]By (A), there exist some $a$ and $b$ in $S$ such that $\gcd(a,b) = 1$, that is, at least one of $a$ or $b$ cannot be divisible by $p$. Denote such an element of $S$ by $\alpha$: thus, $\alpha\not\equiv 0\pmod{p}$. Similarly, by (A), $\gcd(a - 2, b - 2) = 1$, so $p$ cannot divide both $a - 2$ and $b - 2$. Thus, there is an element of $S$, call it $\beta$, such that $\beta\not\equiv 2\pmod{p}$. By $(*)$, $\alpha\equiv 2\pmod{p}$ and $\beta\equiv 0\pmod{p}$. By (B), $\beta^2 - \alpha\in S$. Taking $c = \beta^2 - \alpha$ in $(*)$ yields either $-2\equiv 0\pmod{p}$ or $-2\equiv 2\pmod{p}$, so $p = 2$. Now $(*)$ says that all elements of $S$ are even, contradicting (A). Hence our assumption is false and $S$ is the set of all integers.

 


	\item (\textit{USAMO 2001}) Each point in the plane is assigned a real number such that, for any triangle, the number at the center of its inscribed circle is equal to the arithmetic mean of the three numbers at its vertices. Prove that all points in the plane are assigned the same number.


 



\textbf{Solution}

We label each upper case point with the corresponding lower case letter as its assigned number. The key step is the following lemma.

 \textbf{Lemma}: If $ABCD$ is an isosceles trapezoid, then $a + c = b + d$.

 \textit{Proof}: Assume without loss of generality that $BC\parallel AD$, and that rays $AB$ and $DC$ meet at $P$. Let $I$ be the incenter of triangle $PAC$, and let line $l$ bisect $\angle APD$. Then $I$ is on $l$, so reflecting everything across line $l$ shows that $I$ is also the incenter of triangle $PDB$. Therefore,
 \[\frac{p + a + c}{3} = i = \frac{p + b + d}{3}.\]Hence $a + c = b + d$, as desired.$\blacksquare$

 
                    \begin{figure}[H]
                        \centering
                        \includegraphics[width=0.5\textwidth]{Solution/USAMO_I_2001_6_Solution_Figure_1.png}
                    \end{figure}
                    

 For any two distinct points $A_1$ and $A_2$ in the plane, we construct a regular pentagon $A_1A_2A_3A_4A_5$. Applying the lemma to isosceles trapezoids $A_1A_3A_4A_5$ and $A_2A_3A_4A_5$ yields
 \[a_1 + a_4 = a_3 + a_5\quad\text{and}\quad a_2 + a_4 = a_3 + a_5.\]Hence $a_1 = a_2$. Since $A_1$ and $A_2$ were arbitrary, all points in the plane are assigned the same number.

 


\end{enumerate}